\chapter*{写在前面}
\addcontentsline{toc}{chapter}{写在前面}

我们每天都说自己“知道”：知道明天是否下雨，知道朋友的意图，知道某条新闻不可信，也知道二加二等于四。但这些判断依赖的根据并不相同。天气预报依赖模型和记录，关于朋友的判断依赖证言与交往，数学判断似乎不需要观察。把它们统称为知识，是否掩盖了重要差异？

本书不把认识论写成一串学说和人名，而把它组织成一组仍然开放的问题。每一章从案例开始，继而澄清概念、重构论证、检验反例，最后要求读者亲自作出判断。阅读目标不是记住哪位哲学家“主张什么”，而是学会辨认：一个结论依赖哪些前提，反例究竟攻击了哪里，理论修改又付出了什么代价。

“认识论”在本书中指对知识、确证、理性、理解及认识实践的总体研究；“知识论”只在较窄意义上指知识概念及其条件的分析。英文术语在首次出现时标出，随后使用中文。直接引语注明页码，转述采用作者—年份引证。书中比较材料不是要证明古代中国思想预先提出了某种当代理论，而是检验不同概念框架是否面对相似问题，以及比较在哪些地方会失效。

建议先读案例，在纸上写下直觉，再阅读论证。完成章末练习后才查阅书末提示。哲学自学最常见的误区，是把“读懂”误认为“同意”；真正的理解表现为能够用自己的语言重构最强版本的论证，并说明自己究竟拒绝哪一个前提。

本书面向已修过哲学导论、逻辑学或具备同等阅读能力的高年级本科生，也为硕士阶段的专题课和资格考试复习提供入口。每章末的“分层阅读地图”不是无差别书单：第一层用于补足共同基础，第二层呈现该问题的主要分歧，第三层进入专门研究，第四层连接手册、跨学科材料或可继续写作的议题。初次学习每章选第一层两项和第二层一项即可；研究生研讨应至少覆盖两个相互反对的立场，并从第三、四层追踪它们各自的后续争论。书末参考文献是九章阅读地图的合并结果，每一条均有明确章节归属。

\section*{全书路线}

全书以三部、九个大问题展开。第一部“知识、怀疑与理由”从知识分析、怀疑与认识运气推进到理由和确证结构；第二部“认识来源、社会系统与理性模型”讨论经验、记忆、证言、专家、制度和概率模型怎样共同支持判断与行动；第三部“认识价值、技术环境与比较方法”追问知识、理解与智慧的价值，分析数字环境和人工智能，并用比较认识论重新检查全书采用的问题框架。三部不是彼此封闭的领域：每一部都会把前一部的个人判断问题扩展到来源、制度、价值和方法。每章只保留若干承担完整论证的大节，段内路标用于提示转换，不再把每个细分观点单独切成短小节。
